\section{Álgebras y anillos de polinomios: construcción y propiedades universales}\label{sec:anillos-algebras}

\subsection{Álgebras, morfismos y polinomios formales}\label{tema:13-3-1b}\label{tema:13-3-1d}\label{tema:11-2-1}
El \emph{centro multiplicativo} del anillo es
\[
Z(S)=\set{s\in S\given st=ts\text{ para todo }t\in S}.
\]
\begin{defi}
Sea $R$ conmutativo. Una \emph{$R$-álgebra} es un anillo $S$ junto con un homomorfismo $\alpha:R\to S$ cuya imagen está contenida en $Z(S)$.
\end{defi}
No se exige que $S$ sea conmutativo.

En la categoría $R\text{-}\mathbf{Alg}$, un morfismo de $R$-álgebras $(S,\alpha)\to(T,\beta)$ es un homomorfismo de anillos $f:S\to T$ con $f\alpha=\beta$.

\begin{defi}
Un \emph{polinomio formal} con coeficientes en $R$ es una expresión
\[
f(x)=\sum_{i\geq0}a_ix^i,
\qquad a_i\in R,\quad a_i=0\text{ para todo }i\text{ suficientemente grande}.
\]
Dos polinomios son iguales si coinciden todos sus coeficientes. Su conjunto se denota $R[x]$.
\end{defi}
Como grupo aditivo, $R[x]$ es la suma directa de copias de $(R,+)$ indexadas por los enteros no negativos: cada polinomio es una sucesión de soporte finito. La suma y el producto se definen por
\[
\sum_i a_ix^i+\sum_i b_ix^i=\sum_i(a_i+b_i)x^i,
\qquad
\left(\sum_i a_ix^i\right)\left(\sum_j b_jx^j\right)
=\sum_{k\geq0}\left(\sum_{i+j=k}a_ib_j\right)x^k.
\]
En particular, los primeros coeficientes del producto son
\[
a_0b_0+(a_0b_1+a_1b_0)x
 +(a_0b_2+a_1b_1+a_2b_0)x^2
 +(a_0b_3+a_1b_2+a_2b_1+a_3b_0)x^3+\cdots.
\]
Sólo aparecen finitamente muchos términos no nulos. La unidad es el polinomio constante $1_R$, y $x$ conmuta con los coeficientes aun cuando $R$ no sea conmutativo. La asociatividad del producto resulta de comparar el coeficiente de $x^k$ en ambos agrupamientos: es la suma finita de $a_ib_jc_\ell$ para $i+j+\ell=k$. Las distributivas se comprueban coeficiente a coeficiente. Así, $R[x]$ es un anillo, conmutativo cuando lo es $R$.

Se pueden añadir varias indeterminadas iterando la construcción:
\[
R[x,y,z]=R[x][y][z].
\]
Las nuevas indeterminadas conmutan con todos los coeficientes anteriores y, por tanto, entre sí. La igualdad de polinomios es formal, no una igualdad de sus posibles valores como funciones.

\subsection{Propiedades universales de los anillos de polinomios}\label{tema:11-2-2b}\label{tema:14-3-2a}\label{tema:14-3-2c}
Sea $A=\{a_1,\ldots,a_n\}$. Consideremos la categoría $\mathcal R_A$ cuyos objetos son funciones $j:A\to R$ con $R$ un anillo \emph{conmutativo}. Un morfismo $(j_1,R_1)\to(j_2,R_2)$ es un homomorfismo $\varphi$ que hace conmutar
\[
\begin{tikzcd}[column sep=large]
R_1\arrow[r,"\varphi"]&R_2\\
A\arrow[u,"j_1"]\arrow[ur,"j_2"']&
\end{tikzcd}
\]
\begin{prop}\label{prop:C11-polinomio-universal}
El objeto $(i,\mathbb Z[x_1,\ldots,x_n])$, con $i(a_k)=x_k$, es inicial en $\mathcal R_A$.
\end{prop}
\begin{proof}
Para un objeto $(j,R)$, la conmutatividad exige $\varphi(x_k)=j(a_k)$. El único homomorfismo posible sobre los coeficientes es $\iota_R$, por la \cref{prop:C11-inicial}. Por tanto, todo monomio y toda suma finita están forzados:
\[
\varphi\left(\sum_{\boldsymbol\nu}m_{\boldsymbol\nu}
 x_1^{\nu_1}\cdots x_n^{\nu_n}\right)
=\sum_{\boldsymbol\nu}\iota_R(m_{\boldsymbol\nu})
 j(a_1)^{\nu_1}\cdots j(a_n)^{\nu_n}.
\]
Esto prueba unicidad. La fórmula define una función que preserva suma, unidad y producto: al multiplicar, la distributividad reúne los términos de exponentes cuya suma es la misma, y la conmutatividad de $R$ permite ordenar los factores. Además, envía cada $x_k$ a $j(a_k)$. Es, pues, el mediador requerido:
\[
\begin{tikzcd}[column sep=large]
\mathbb Z[x_1,\ldots,x_n]\arrow[r,destacado,"\exists!\,\varphi"]&R\\
A\arrow[u,"i"]\arrow[ur,"j"']&
\end{tikzcd}
\qedhere
\]
\end{proof}

Supongamos $R$ conmutativo y $A=\{1,\ldots,n\}$. Ponemos $R[A]=R[x_1,\ldots,x_n]$ y $j(i)=x_i$.
\begin{prop}\label{prop:C14-algebra-libre}
Para toda $R$-álgebra \emph{conmutativa} $(S,\alpha)$ y función $f:A\to S$, existe un único homomorfismo de $R$-álgebras $\varphi:R[A]\to S$ con $\varphi(x_i)=f(i)$.
\end{prop}
\begin{proof}
Extendemos el homomorfismo sucesivamente. Primero, la fórmula
\[
\varphi_1\left(\sum_k r_kx_1^k\right)=\sum_k\alpha(r_k)f(1)^k
\]
extiende $\alpha$ a $R[x_1]$: preserva suma, producto y unidad porque $f(1)$ conmuta con los coeficientes. Es única porque fija los coeficientes y envía $x_1$ a $f(1)$. La compatibilidad con el producto se comprueba expandiendo:
\[
\varphi_1(r_0+r_1x_1^{k_1})=\alpha(r_0)+\alpha(r_1)f(1)^{k_1}.
\]
Después extendemos a $R[x_1][x_2]$, enviando $x_2$ a $f(2)$; este elemento conmuta con toda la imagen anterior porque $S$ es conmutativo. Continuamos hasta $x_n$. Cada extensión es única y la fórmula resultante es
\[
\varphi\left(\sum_{\boldsymbol\nu}r_{\boldsymbol\nu}
x_1^{\nu_1}\cdots x_n^{\nu_n}\right)
=\sum_{\boldsymbol\nu}\alpha(r_{\boldsymbol\nu})
f(1)^{\nu_1}\cdots f(n)^{\nu_n}.
\]
Su restricción a las constantes es $\alpha$, así que es morfismo de $R$-álgebras. Los triángulos sucesivos se reúnen en el siguiente, conservando las imágenes de todos los generadores:
\[
\begin{tikzcd}[column sep=large,row sep=large]
R[A]\arrow[r,dashed,destacado,"\exists!\,\varphi"]&S\\
A\arrow[u,"j"]\arrow[ur,"f"']&
\end{tikzcd}
\qedhere
\]
\end{proof}

\begin{obs}
El homomorfismo estructural $\alpha:R\to S$ no tiene por qué ser inyectivo: la condición central es $\alpha(R)\subseteq Z(S)$. Para evaluar varias indeterminadas, sus imágenes deben conmutar entre sí. Por eso la propiedad universal de $R[A]$ se formula aquí para álgebras conmutativas.
\end{obs}
